\emph{The algebra.} In $\Cl(1,2)$, $e_{1}^{2} = 1$, $e_{2}^{2} = e_{3}^{2} = -1$: $I = e_{1}e_{2}e_{3}$ commutes with each generator (three generators, each anticommuting with the two others) and $I^{2} = -e_{1}^{2}e_{2}^{2}e_{3}^{2} = -1$. Then $(Ie_{2})^{2} = I^{2}e_{2}^{2} = 1$, $(Ie_{3})^{2} = 1$, and $e_{1}$, $Ie_{2}$, $Ie_{3}$ anticommute. $Ie_{3} = e_{1}e_{2}e_{3}^{2} = -e_{1}e_{2}$, so $P_{1} = vw = \tfrac12(1 - e_{1}e_{2}) = \tfrac12(1 + Ie_{3})$ (the form of $vw$ is in the stored proof, \claimref{synthesis.set-and-clear-of-a-bit-generate-the-hyperbolic-plane}). For unit combinations $m$, $m'$ of $e_{1}, Ie_{2}, Ie_{3}$, $\tfrac12(1 + m)$ is an idempotent fixed by the conjugation and $2\br{\tfrac12(1 + m)\tfrac12(1 + m')}_{0} = \tfrac12(1 + m \cdot m')$ lies in $[0,1]$, where $m \cdot m'$ is the scalar part of $mm'$: $\langle\,\cdot\,\rangle$ is an expectation and its values on such idempotents are probabilities. The conjugation sends $I$ to $(-e_{3})(-e_{2})e_{1} = -I$.

\emph{(1).} $J$ is a combination of the idempotents $P_{1}$, $P_{2}$ with $P_{1}P_{2} = 0$, $P_{1} + P_{2} = 1$, and $I$ is central with $I^{2} = -1$, so $\exp(-I\alpha J) = e^{-I\alpha c_{1}}P_{1} + e^{-I\alpha c_{2}}P_{2}$ with $e^{-I\beta} = \cos\beta - I\sin\beta$; its conjugate is $e^{I\alpha c_{1}}P_{1} + e^{I\alpha c_{2}}P_{2}$, its inverse. With no free part clause (2) of Theorem \claimref{synthesis.field-slaved-to-an-operator-source-does-not-commute-at-spacelike-separation} gives $\phi_{n}(t) = gJ\,S(n,t)$, an element of the matter algebra, so (D) gives $V = \exp(-I\lambda gS\,J)$; the matter's own coupling to this field is a function of $J$ and is a unit in the same way. The stored clause (2) is the statement that $[\phi_{n}(t), O] \neq 0$ at spacelike events; so here is a unitary model in which the condition fails.

\emph{(2).} $V = e^{-I\lambda gSc_{1}}\bigl(P_{1} + e^{I\vartheta}P_{2}\bigr)$. With $e_{1} = v + w$, $v = P_{1}vP_{2}$, $w = P_{2}wP_{1}$: $V^{-1}vV = e^{I\vartheta}v$ and $V^{-1}wV = e^{-I\vartheta}w$, so $V^{-1}e_{1}V = \cos\vartheta\,(v + w) + I\sin\vartheta\,(v - w) = e_{1}\cos\vartheta + Ie_{2}\sin\vartheta$, while $V$ commutes with $P_{1}$, $P_{2}$. From the stored proof, $O = p^{2}P_{1} + q^{2}P_{2} + pq\,e_{1}$, so $V^{-1}OV - O = pq\,\bigl(e_{1}(\cos\vartheta - 1) + Ie_{2}\sin\vartheta\bigr)$, and $\langle e_{1}\rangle = a$, $\langle Ie_{2}\rangle = b$ in the state $\rho$. A device placed not later than the detection acts before the recombination, so the detection probability is $\langle V^{-1}OV\rangle$. To first order $V^{-1}OV = O + I\lambda[\phi_{n}(t), O]$, and $[\phi_{n}(t), O] = gS\,[J, O] = gS\,(c_{1} - c_{2})pq\,e_{2}$ by the stored clause (1); $\langle Ie_{2}\rangle = b$. For the bound: $a(\cos\vartheta - 1) + b\sin\vartheta = 2\sin(\vartheta/2)\,\bigl(-a\sin(\vartheta/2) + b\cos(\vartheta/2)\bigr)$, of absolute value at most $2|\sin(\vartheta/2)|$ for $a^{2} + b^{2} = 1$, and $pq \le \tfrac12$ for $p^{2} + q^{2} = 1$.

\emph{(3).} The count of the events with $S \neq 0$ is the stored one ($172$ of $502$); those with $t \le 8$ are counted here ($32$ of $252$). For an event later than the detection the device acts after it and cannot change it; there the nonzero commutator says that the expectation of $\phi_{n}(t)$ depends on whether $O$ was measured (a measurement of $O$ whose outcome is not kept removes the part of the state that does not commute with $O$, and $J$ does not commute with $O$).

\emph{(4).} Clause (4) of the stored result: with the free part, $[\phi_{n}(t), O] = -g[J, O]\sum_{t'}(G + \Delta)$ to first order, zero at the $502$ spacelike events; the coefficient of (2) is its expectation with $I$. (The stored proof notes that the statement is exact, the kicks being displacement operators; not used.)

\emph{(5).} Under the stored coupling of Proposition \claimref{field-theory.operator-newtonian-coupling} the branch with body $A$ on path $i$ has energy $-Gm_{A}m_{B}/r_{i}$ with a mass $m_{B}$ at distance $r_{i}$, so the phase between the branches grows at the rate $Gm_{A}m_{B}|1/r_{1} - 1/r_{2}|/\hbar$. Reading the field of $A$ at the device's position as a function of $A$'s path number at the retarded time, which for a superposition older than $D/c$ is the same function, and coupling the device to it by (D), gives this rate from the moment the device is in place; a device placed within the time $D/c$ before the detection is at events spacelike to it. This is the argued step: it uses the extension of the element reading made in the stored result's Reading, and (D). The numbers: $G/(\hbar c) = 2.111 \times 10^{15}\ \mathrm{kg}^{-2}$ (from $G = 6.674 \times 10^{-11}$, $\hbar = 1.0546 \times 10^{-34}$, $c = 2.998 \times 10^{8}$ in SI units), the inverse square of the Planck mass $2.18 \times 10^{-8}$ kg; $D|1/r_{1} - 1/r_{2}| = Dd/(D^{2} - d^{2}/4)$, which is $180/299$ for $d = 250$, $D = 450$ and $2.5 \times 10^{-4}$ for $d = 250\,\mu$m, $D = 1$ m. Then $\Phi_{w} = 2.111 \times 10^{15} \times 10^{-28} \times 0.602 = 1.27 \times 10^{-13}$; the change of probability is at most $|\sin(\vartheta/2)| \le \Phi_{w}/2$; the number of runs with $\tfrac12\Phi_{w}\sqrt{N} = 1$ is $4/\Phi_{w}^{2} = 2.5 \times 10^{26}$. For one kilogram at one metre, $2.111 \times 10^{15} \times 10^{-14} \times 2.5 \times 10^{-4} = 5.3 \times 10^{-3}$.

\emph{(6).} In Planck units a mass of $10^{3}$ kg is $M = 10^{3}/(2.18 \times 10^{-8}) = 4.6 \times 10^{10}$, and $r = 0.1$ m is $6.2 \times 10^{33}$, so $2M/r^{3} = 4 \times 10^{-91}$. The level is one at $r = (2M)^{1/3} = 4.5 \times 10^{3}$ Planck lengths, $7 \times 10^{-32}$ m, less than $2M = 9.2 \times 10^{10}$ Planck lengths, $1.5 \times 10^{-24}$ m. Premises \claimref{singularity.c1-curvature-cutoff} and \claimref{singularity.c2-cell-density} speak of that surface for a hole.

\emph{(S).} Each item is a reading of the cited entity's statement at the pinned version: Theorem \claimref{synthesis.locality-of-the-attached-free-fields-fixes-their-statistics} (its clause 1 lists the scalar, the spinor and the photon); Fact \claimref{synthesis.imported-dimock-dirac-car-algebra-on-a-gravitational-field} (its clause (a)); Theorem \claimref{synthesis.interacting-fields-of-causal-perturbation-theory-keep-the-spin-statistics-connection-at-even-ghost-number} (its clause (i) lists the interaction monomials); Premise \claimref{synthesis.c-loc-local-horizon-premise} (clause (Loc): a density fixed by flat tangent data, calibrated with no mode count). A text search of the store for commutation at spacelike separation, signalling and a least length found no further statement. Items (d) and (e) follow from the form of the stored coupling and from (1), (2).
